% Lösungen Einheit 3 – Nr. 33–43
\einheitenkopf[le3][3 Kreisteile]{Einheit 3 von 3 · Kreisteile}
\erg{33}{1. Halbkreis \quad 2. Viertelkreis \quad 3. Dreiviertelkreis \quad 4. Achtelkreis \quad 5. Drittelkreis}
\erg{34}{a) $\tfrac{1}{4}$ \quad b) $\tfrac{1}{2}$ \quad c) $\tfrac{1}{6}$ \quad d) $\tfrac{1}{3}$ \quad e) 10 \% \quad f) 27,8 \% \quad g) grau sind $360^\circ - 125^\circ = 235^\circ$; $235 : 360 \approx 65{,}3\,\%$}
\erg{35}{a) 90° \quad b) 240° \quad c) $0{,}3 \cdot 360^\circ = 108^\circ$}
\erg{36}{a) 6,3 cm \quad b) 9,4 cm \quad c) 18,8 cm \quad d) $r = 2{,}5\,\text{cm}$: 10,9 cm}
\erg{37}{a) 12,6 cm² \quad b) 37,7 cm² \quad c) 98,2 cm² \quad d) $r = 4{,}5\,\text{m}$: 53,0 m²}
\erg{38}{a) $6{,}28 + 8 \approx 14{,}3\,\text{cm}$ \quad b) $12{,}57 + 24 \approx 36{,}6\,\text{cm}$ \quad c) $13{,}09 + 10 \approx 23{,}1\,\text{cm}$}
\erg{39}{a) $\pi \cdot (36 - 16) \approx 62{,}8\,\text{cm}^2$ \quad b) $\pi \cdot (15^2 - 10^2) \approx 392{,}7\,\text{cm}^2$ \quad c) innerer Radius 6 cm: $\pi \cdot (81 - 36) \approx 141{,}4\,\text{cm}^2$}
\erg{40}{a) durch 180 statt durch 360 geteilt; richtig: $40 : 360 = \tfrac{1}{9}$ \quad b) die beiden Radien fehlen; richtig: $9{,}42 + 2 \cdot 6 \approx 21{,}4\,\text{cm}$ \quad c) mit dem weißen Restwinkel gerechnet; grau sind $250^\circ$: $250 : 360 \approx 69{,}4\,\%$}
\erg{41}{a) Ein Durchmesser teilt den Kreis in zwei gleiche Hälften; der Vollwinkel am Mittelpunkt ist $360^\circ$, die Hälfte davon $180^\circ$ (gestreckter Winkel). \quad b) Gleicher Anteil, $\tfrac{1}{6}$, weil nur der Winkel den Anteil bestimmt; die Flächen sind verschieden, weil $\tfrac{1}{6}$ eines größeren Kreises größer ist.}
\erg{42}{a) $\tfrac{150}{360} \cdot \pi \cdot 7^2 \approx 64{,}1\,\text{m}^2$ \quad b) $\tfrac{150}{360} \cdot 2 \cdot \pi \cdot 7 \approx 18{,}3\,\text{m}$ \quad c) $80 : 153{,}94 \approx 0{,}52$; $0{,}52 \cdot 360^\circ \approx 187^\circ$}
\erg{43}{a) Rad 45 \%, Bus 30 \%, zu Fuß 25 \% \quad b) Rad 162°, Bus 108°, zu Fuß 90° \quad c) Kreisdiagramm:}

\kreisdiagramm[1.2]{Rad/162, Bus/108, zu Fuß/90}
